% !PRINTED 650
\weberpage{650}
%<<CONT
\weq{7}{\begin{aligned}k_0+k_1+k_2+k_3 &= 4\,C,\\ k_0-k_1+k_2-k_3 &= 4\,B\sqrt{c},\end{aligned}}
also
\weq{8}{\begin{aligned}k_0+k_2 &= 2\,(C+B\sqrt{c}),\\ k_1+k_3 &= 2\,(C-B\sqrt{c}),\end{aligned}}
wodurch nach (5) folgt:
\weq{9}{\begin{aligned}k_0 &= C+B\sqrt{c}+\sqrt{b+a\sqrt{c}}\\ k_1 &= C-B\sqrt{c}+\sqrt{b-a\sqrt{c}}\\ k_2 &= C+B\sqrt{c}-\sqrt{b+a\sqrt{c}}\\ k_3 &= C-B\sqrt{c}-\sqrt{b-a\sqrt{c}}.\end{aligned}}
Es ist auch umgekehrt leicht zu zeigen, dass diese Grössen
die Wurzeln einer biquadratischen cyklischen Gleichung sind.
Denn setzen wir zur Abkürzung
\weq{10}{r = \sqrt{c}, \quad \varrho = \sqrt{b+a\sqrt{c}}, \quad \varrho' = \sqrt{b-a\sqrt{c}},}
also nach (6)
\[ r = \varrho\varrho', \quad \varrho^2 = b+ar, \quad \varrho'^2 = b-ar, \]
so können wir jede rationale Function der $k_0, k_1, k_2, k_3$ als lineare
Function mit rationalen Coëfficienten von den sechs Radicalen
\[ 1, \quad r, \quad \varrho, \quad \varrho', \quad r\varrho, \quad r\varrho' \]
darstellen. Die cyklische Permutation $(k_0, k_1, k_2, k_3)$ entspricht
der Vertauschung
\[
\begin{array}{lll}
r, & \varrho, & \varrho',\\
-r, & \varrho', & -\varrho,
\end{array}
\]
und wenn man diese Vertauschung wiederholt, so sieht man,
dass eine cyklische Function der $k$ sich nicht ändern kann,
wenn in der linearen Darstellung durch diese sechs Grössen
folgende Vertauschungen gemacht werden:
\[
\begin{array}{llllll}
1, & r, & \varrho, & \varrho', & r\varrho, & r\varrho'\\
1, & -r, & \varrho', & -\varrho, & -r\varrho', & r\varrho\\
1, & r, & -\varrho, & -\varrho', & -r\varrho, & -r\varrho'\\
1, & -r, & -\varrho', & \varrho, & r\varrho', & -r\varrho,
\end{array}
\]
und wenn man die vier so sich ergebenden Ausdrücke addirt,
so erhält man für die cyklische Function der $k$ einen rationalen
Ausdruck.

Will man die biquadratische Gleichung bilden, deren Wurzeln
die $k$ sind, so führt man am besten $k - C = x$ als Unbekannte
%CONT>>
